% source: RKF/theorum/46_first_visible_jet_seam_quotient_theorem.md
\hypertarget{first-visible-jet-seam-quotient-closure-theorem}{%
\chapter{First-Visible-Jet Seam Quotient Closure Theorem}\label{first-visible-jet-seam-quotient-closure-theorem}}

\hypertarget{purpose-and-claim-boundary}{%
\section{1. Purpose and claim boundary}\label{purpose-and-claim-boundary}}

This theorem gives a rigorous finite-jet resolution of a class of apparent
\texttt{0/0} indeterminate forms. It does \textbf{not} redefine ordinary algebraic division
by zero. The raw expressions

\begin{verbatim}
1/0
0/0
\end{verbatim}

remain invalid as algebraic quotients.

The object defined here is instead a \textbf{seam quotient of two functions} along a
declared one-parameter seam. Its value is determined, when possible, by the
first nonvanishing relative jet together with explicit remainder control.

The construction consumes the current capstone chain:

\begin{verbatim}
Theorem 42  typed derivative / response tower
Theorem 43  first-visible jet and explicit remainder
Theorem 44  body-tail transfer under refinement
Theorem 45  arithmetic radius + analytic tail enclosure
\end{verbatim}

The motivating Dot-in-Circle seam principle is therefore narrowed to a theorem:
indeterminate endpoint values may be resolved only when the declared seam
carries enough nonvanishing local structure to determine a unique quotient.

\begin{center}\rule{0.5\linewidth}{0.5pt}\end{center}

\hypertarget{declared-seam-and-scalar-carrier}{%
\section{2. Declared seam and scalar carrier}\label{declared-seam-and-scalar-carrier}}

Let \texttt{K} be either \texttt{R} or \texttt{C}. Let

\[
S:(-\varepsilon,\varepsilon)\to X
\]

be a declared local seam with distinguished endpoint \texttt{S(0)}. Let

\[
A,B:(-\varepsilon,\varepsilon)\to K
\]

be scalar observables evaluated along this seam. We write simply \texttt{A(t)} and
\texttt{B(t)}.

Assume

\[
A(0)=B(0)=0.
\tag{2.1}
\]

For an integer \texttt{N\ \textgreater{}=\ 0}, suppose the two observables admit finite jet
expansions

\[
A(t)=\sum_{k=0}^{N} a_k t^k + R_A^{(N)}(t),
\tag{2.2}
\]

\[
B(t)=\sum_{k=0}^{N} b_k t^k + R_B^{(N)}(t),
\tag{2.3}
\]

where the coefficient normalization may absorb factorials. In a derivative
normalization one has \texttt{a\_k=A\^{}\{(k)\}(0)/k!} and similarly for \texttt{b\_k}.

Define the finite first-visible orders

\[
r_A=\min\{k\le N:a_k\ne0\},
\qquad
r_B=\min\{k\le N:b_k\ne0\},
\tag{2.4}
\]

whenever the sets are nonempty.

If every coefficient through the available depth vanishes, the corresponding
finite first-visible order is \textbf{unresolved}, not infinity by decree.

\begin{center}\rule{0.5\linewidth}{0.5pt}\end{center}

\hypertarget{first-visible-jet-classification}{%
\section{3. First-visible-jet classification}\label{first-visible-jet-classification}}

\hypertarget{theorem-3.1-finite-first-visible-jet-quotient-classification}{%
\subsection{Theorem 3.1 (Finite first-visible-jet quotient classification)}\label{theorem-3.1-finite-first-visible-jet-quotient-classification}}

Assume finite first-visible orders \texttt{r\_A} and \texttt{r\_B} exist and that the remainders
are little-o at the corresponding leading orders:

\[
R_A^{(N)}(t)=o(t^{r_A}),
\qquad
R_B^{(N)}(t)=o(t^{r_B}).
\tag{3.1}
\]

Then along the declared seam:

\begin{enumerate}
\def\labelenumi{\arabic{enumi}.}
\tightlist
\item
  if \texttt{r\_A\ \textgreater{}\ r\_B},
\end{enumerate}

\[
\boxed{\lim_{t\to0}\frac{A(t)}{B(t)}=0;}
\tag{3.2}
\]

\begin{enumerate}
\def\labelenumi{\arabic{enumi}.}
\setcounter{enumi}{1}
\tightlist
\item
  if \texttt{r\_A\ =\ r\_B\ =\ r},
\end{enumerate}

\[
\boxed{\lim_{t\to0}\frac{A(t)}{B(t)}=\frac{a_r}{b_r};}
\tag{3.3}
\]

\begin{enumerate}
\def\labelenumi{\arabic{enumi}.}
\setcounter{enumi}{2}
\tightlist
\item
  if \texttt{r\_A\ \textless{}\ r\_B}, then
\end{enumerate}

\[
\boxed{\left|\frac{A(t)}{B(t)}\right|\to\infty,}
\tag{3.4}
\]

so no finite seam quotient exists.

\hypertarget{proof}{%
\subsubsection{Proof}\label{proof}}

Factor the leading powers:

\[
A(t)=t^{r_A}(a_{r_A}+o(1)),
\qquad
B(t)=t^{r_B}(b_{r_B}+o(1)).
\]

Since \texttt{b\_\{r\_B\}\ !=\ 0}, the second parenthesis is nonzero for sufficiently small
punctured \texttt{t}. Hence

\[
\frac{A(t)}{B(t)}
=t^{r_A-r_B}
\frac{a_{r_A}+o(1)}{b_{r_B}+o(1)}.
\]

The three cases follow from the sign of \texttt{r\_A-r\_B}. \texttt{square}

\begin{center}\rule{0.5\linewidth}{0.5pt}\end{center}

\hypertarget{dot-in-circle-seam-quotient}{%
\section{4. Dot-in-Circle seam quotient}\label{dot-in-circle-seam-quotient}}

\hypertarget{definition-4.1-certified-dot-in-circle-quotient}{%
\subsection{Definition 4.1 (Certified Dot-in-Circle quotient)}\label{definition-4.1-certified-dot-in-circle-quotient}}

When Theorem 3.1 gives a finite limit, define

\[
\boxed{A\odot_S B:=\lim_{t\to0}\frac{A(t)}{B(t)}.}
\tag{4.1}
\]

The subscript \texttt{S} is part of the type. Different seams need not produce the same
value unless a seam-invariance theorem is separately proved.

In the equal-order case,

\[
\boxed{A\odot_S B=\frac{a_r}{b_r}.}
\tag{4.2}
\]

This is not the algebraic statement \texttt{0/0=a\_r/b\_r}. It is the limit statement
that two vanishing observables retain a first nonvanishing relative jet along a
declared seam.

\hypertarget{example-4.2}{%
\subsection{Example 4.2}\label{example-4.2}}

Let

\[
A(t)=t\cos\alpha,
\qquad
B(t)=t\sin\alpha.
\]

If \texttt{sin(alpha)\ !=\ 0}, then \texttt{r\_A=r\_B=1} and

\[
\boxed{A\odot_S B=\cot\alpha.}
\tag{4.3}
\]

\begin{center}\rule{0.5\linewidth}{0.5pt}\end{center}

\hypertarget{regular-seam-reparameterization-invariance}{%
\section{5. Regular seam-reparameterization invariance}\label{regular-seam-reparameterization-invariance}}

A seam may be described by more than one regular local parameter. Let

\[
t=\phi(s),
\qquad
\phi(0)=0,
\qquad
\phi'(0)=c\ne0.
\tag{5.1}
\]

\hypertarget{theorem-5.1-regular-parameter-invariance}{%
\subsection{Theorem 5.1 (Regular parameter invariance)}\label{theorem-5.1-regular-parameter-invariance}}

Suppose \texttt{A} and \texttt{B} have equal first-visible order \texttt{r}. Under the regular
reparameterization (5.1), their leading coefficients become

\[
\widetilde a_r=a_r c^r,
\qquad
\widetilde b_r=b_r c^r.
\tag{5.2}
\]

Therefore

\[
\boxed{\frac{\widetilde a_r}{\widetilde b_r}=\frac{a_r}{b_r}.}
\tag{5.3}
\]

The finite Dot-in-Circle seam quotient is invariant under regular local
reparameterization of the same seam.

\hypertarget{proof-1}{%
\subsubsection{Proof}\label{proof-1}}

Since \texttt{phi(s)=cs+o(s)}, one has

\[
A(\phi(s))=a_r c^r s^r+o(s^r),
\qquad
B(\phi(s))=b_r c^r s^r+o(s^r).
\]

Cancel the common nonzero factor \texttt{c\^{}r}. \texttt{square}

The theorem does not identify genuinely different geometric paths. It only
removes arbitrary regular changes of local seam coordinate.

\begin{center}\rule{0.5\linewidth}{0.5pt}\end{center}

\hypertarget{quantitative-denominator-separation}{%
\section{6. Quantitative denominator separation}\label{quantitative-denominator-separation}}

The limit theorem becomes proof-bearing numerically only after the reduced
denominator is separated from zero.

Assume the common first-visible order is \texttt{r} and define the reduced functions

\[
\widehat A(t)=\frac{A(t)}{t^r},
\qquad
\widehat B(t)=\frac{B(t)}{t^r}
\quad(t\ne0).
\tag{6.1}
\]

Suppose on a certified punctured seam neighborhood,

\[
|\widehat A(t)-a|\le\delta_A,
\qquad
|\widehat B(t)-b|\le\delta_B,
\tag{6.2}
\]

with

\[
\boxed{|b|>\delta_B.}
\tag{6.3}
\]

\hypertarget{theorem-6.1-certified-quotient-enclosure}{%
\subsection{Theorem 6.1 (Certified quotient enclosure)}\label{theorem-6.1-certified-quotient-enclosure}}

Under (6.2)--(6.3), \texttt{widehat\ B(t)} is nonzero throughout the certified
neighborhood and

\[
\boxed{
\left|
\frac{A(t)}{B(t)}-\frac{a}{b}
\right|
\le
\frac{|b|\delta_A+|a|\delta_B}
{|b|\,(|b|-\delta_B)}.
}
\tag{6.4}
\]

\hypertarget{proof-2}{%
\subsubsection{Proof}\label{proof-2}}

Write

\[
\widehat A=a+e_A,
\qquad
\widehat B=b+e_B,
\]

with \texttt{\textbar{}e\_A\textbar{}\textless{}=delta\_A}, \texttt{\textbar{}e\_B\textbar{}\textless{}=delta\_B}. Condition (6.3) gives

\[
|b+e_B|\ge |b|-\delta_B>0.
\]

Then

\[
\frac{a+e_A}{b+e_B}-\frac ab
=\frac{b e_A-a e_B}{b(b+e_B)}.
\]

Taking absolute values and applying the triangle inequality gives (6.4).
\texttt{square}

The quantity

\[
\boxed{
\rho_Q=
\frac{|b|\delta_A+|a|\delta_B}
{|b|\,(|b|-\delta_B)}
}
\tag{6.5}
\]

is the outward quotient radius. It may be consumed by Theorem 45 as a numerical
arithmetic/analytic enclosure only when \texttt{delta\_A} and \texttt{delta\_B} themselves are
proved bounds.

\begin{center}\rule{0.5\linewidth}{0.5pt}\end{center}

\hypertarget{connection-to-theorems-4245}{%
\section{7. Connection to Theorems 42--45}\label{connection-to-theorems-4245}}

The capstone chain supplies the data needed by the quotient theorem.

\hypertarget{theorem-42-typed-tower}{%
\subsection{7.1 Theorem 42: typed tower}\label{theorem-42-typed-tower}}

The derivative/response tower provides the candidate coefficients

\[
a_k,\ b_k.
\]

The first nonzero compatible layer determines the first-visible order.

\hypertarget{theorem-43-first-visibility-and-remainder}{%
\subsection{7.2 Theorem 43: first visibility and remainder}\label{theorem-43-first-visibility-and-remainder}}

The nested observer tower identifies the first order at which the target ceases
to be blind. Its explicit finite-jet remainder gives a lawful route to bounds
of the form \texttt{delta\_A}, \texttt{delta\_B}.

\hypertarget{theorem-44-no-discarded-tail}{%
\subsection{7.3 Theorem 44: no discarded tail}\label{theorem-44-no-discarded-tail}}

Refinement transfers one term from tail to recognized body while preserving

\[
M_r+T_r=F.
\]

Thus a quotient certificate may not erase higher-order uncertainty after the
first visible coefficient is found.

\hypertarget{theorem-45-arithmetic-enclosure}{%
\subsection{7.4 Theorem 45: arithmetic enclosure}\label{theorem-45-arithmetic-enclosure}}

Exact rational or finite-decimal coefficients have zero arithmetic radius.
Approximate coefficients require validated interval/ball enclosures. The
quotient radius (6.5) and any remaining analytic tail are carried outward before
promotion.

\begin{center}\rule{0.5\linewidth}{0.5pt}\end{center}

\hypertarget{flat-function-boundary}{%
\section{8. Flat-function boundary}\label{flat-function-boundary}}

Finite jets do not resolve every \texttt{0/0} limit.

For example,

\[
A(t)=e^{-1/t^2},
\qquad
B(t)=e^{-1/t^2}
\quad(t\ne0),
\]

extended by zero at \texttt{t=0}, have all derivatives zero at the endpoint but

\[
\frac{A(t)}{B(t)}=1
\]

on the punctured seam.

Therefore:

\[
\boxed{
\text{all available jets zero}
\not\Rightarrow
\text{quotient absent};
}
\tag{8.1}
\]

it means only that the finite-jet theorem has insufficient information.

The fail-closed status is

\begin{verbatim}
INCOMPLETE_FLAT_OR_UNRESOLVED
\end{verbatim}

until a different admitted seam invariant or non-jet asymptotic theorem closes
the case.

This boundary is essential. Declaring all flat \texttt{0/0} forms equal would simply
replace indeterminacy with invention.

\begin{center}\rule{0.5\linewidth}{0.5pt}\end{center}

\hypertarget{raw-division-by-zero-boundary}{%
\section{9. Raw division-by-zero boundary}\label{raw-division-by-zero-boundary}}

The theorem distinguishes the following typed cases:

\begin{verbatim}
1 / 0 as algebraic scalar                     INVALID
0 / 0 as algebraic scalar                     INVALID
A(t)/B(t), no declared seam                    INVALID_CONTRACT
seam declared, no finite denominator jet       INCOMPLETE_FLAT_OR_UNRESOLVED
r_A > r_B                                      FINITE_QUOTIENT_ZERO
r_A = r_B and b_r != 0                         FINITE_SEAM_QUOTIENT
r_A < r_B                                      DIVERGENT_NO_FINITE_QUOTIENT
reduced denominator not separated from zero    INCOMPLETE_DENOMINATOR_SEPARATION
\end{verbatim}

No branch converts \texttt{1/0} into a finite certified number.

\begin{center}\rule{0.5\linewidth}{0.5pt}\end{center}

\hypertarget{optional-phase-compatibility}{%
\section{10. Optional phase compatibility}\label{optional-phase-compatibility}}

The Bindu-Lopa source motivates a phase lift on the unit circle. Such a phase
layer may be consumed only \textbf{after} the scalar seam quotient has been resolved
or enclosed.

Suppose a separately admitted phase factor satisfies

\[
\Phi(t)\to\Phi_0,
\qquad
|\Phi_0|=1.
\]

If \texttt{A\ odot\_S\ B=q} exists, then

\[
\boxed{\lim_{t\to0}\frac{A(t)}{B(t)}\Phi(t)=q\Phi_0.}
\tag{10.1}
\]

The phase factor decorates an already resolved quotient; it does not by itself
recover discarded magnitude information or prove existence of the quotient.

\begin{center}\rule{0.5\linewidth}{0.5pt}\end{center}

\hypertarget{negative-controls-required-for-certification}{%
\section{11. Negative controls required for certification}\label{negative-controls-required-for-certification}}

Any implementation claiming this theorem must include at least:

\begin{enumerate}
\def\labelenumi{\arabic{enumi}.}
\tightlist
\item
  raw \texttt{1/0} rejection;
\item
  raw \texttt{0/0} rejection;
\item
  equal-order finite quotient fixture;
\item
  numerator-higher-order zero-limit fixture;
\item
  denominator-higher-order divergence fixture;
\item
  all-zero finite-jet incomplete fixture;
\item
  regular seam-reparameterization invariance fixture;
\item
  denominator-separation failure fixture;
\item
  exact-rational coefficient fixture;
\item
  quotient-enclosure bound fixture.
\end{enumerate}

A numerical implementation must additionally prove or validate the coefficient
and remainder bounds it consumes. Merely labelling a submitted radius
\texttt{validated} does not establish the hypothesis of Theorem 6.1.

\begin{center}\rule{0.5\linewidth}{0.5pt}\end{center}

\hypertarget{claim-boundary}{%
\section{12. Claim boundary}\label{claim-boundary}}

The proved statement is:

\begin{quote}
A pair of vanishing scalar observables along a declared regular seam has a
finite seam quotient whenever their first nonzero jets occur at the same
finite order and the denominator leading coefficient is nonzero; the quotient
equals the ratio of those leading coefficients. Unequal first-visible orders
classify zero versus divergence. Explicit reduced-function remainder bounds
give an outward quotient enclosure once the reduced denominator is separated
from zero.
\end{quote}

The theorem does \textbf{not} claim:

\begin{itemize}
\tightlist
\item
  a universal value for algebraic \texttt{0/0};
\item
  a finite value for \texttt{1/0};
\item
  uniqueness across genuinely different seams without an additional invariance theorem;
\item
  finite-jet resolution of flat functions;
\item
  correctness of an externally asserted arithmetic radius or analytic tail;
\item
  that phase lifting alone resolves magnitude indeterminacy.
\end{itemize}

This theorem is therefore a conditional closure theorem for indeterminate
limits, not a redefinition of field arithmetic.

\begin{flushright}\small\texttt{RKF/theorum/46\_first\_visible\_jet\_seam\_quotient\_theorem.md}\end{flushright}
